\exercisegroup

¶ 1) Let E be a locally convex metrizable space, and \(E_{b}^{\prime}\) its strong dual. If \(E_{b}^{\prime}\) is metrizable, prove that the topology of E can be defined by a single norm (use III, p.~37, exerc. 2 and p.~38, exerc. 5 and also the fact that E is bornological).

\begin{enumerate}
\def\labelenumi{\arabic{enumi})}
\setcounter{enumi}{1}
\tightlist
\item
  An infra-barrelled space is semi-barrelled. A locally convex space is said to be a (DF) space if it is semi-barrelled and if the canonical bornology (III, p.~3, def. 5) has a countable base. Every normed space and every strict inductive limit of a sequence of normed spaces (II, p.~33) is a (DF) space. Every strong dual of a Fréchet space is a (DF) space.
\end{enumerate}

\begin{enumerate}
\def\labelenumi{\alph{enumi})}
\item
  The strong dual of a (DF) space is a Fréchet space.
\item
  Let \(\mathrm{E}\) be a (DF) space and let \((\mathbf{A}_n)\) be an increasing sequence of bounded, convex, balanced and closed subsets of \(\mathrm{E}\) such that every bounded subset of \(\mathrm{E}\) is absorbed by one of the \(\mathbf{A}_n\) .
\end{enumerate}

Let U be the union of the \(A_{n}\) ; show that the closure \(\overline{U}\) of U in E is precisely the set of all \(x \in E\) such that \(\lambda x \in U\) for \(0 \leqslant \lambda < 1\) . (If \(x \notin \lambda U\) for some \(\lambda > 1\) , then for every n, there exists a linear form \(x_{n}' \in E'\) such that \(x_{n}' \in A_{n}^{\circ}\) and \(\langle x, x_{n}' \rangle = \lambda\) , and the sequence \((x_{n}')\) is equicontinuous, hence has a weak limit point.)

\begin{enumerate}
\def\labelenumi{\alph{enumi})}
\setcounter{enumi}{2}
\tightlist
\item
  Show that if a (DF) space is barrelled, it is also bornological (cf.~III, p.~44, exerc. 13, b)). Give examples of (DF) spaces which are not ultrabornological, but are bornological and barrelled (III, p.~46, exerc. 22) and also of (DF) spaces which are not barrelled but are bornological.
\end{enumerate}

\begin{enumerate}
\def\labelenumi{\arabic{enumi})}
\setcounter{enumi}{2}
\tightlist
\item
  Let \(E\) be a locally convex metrizable space, and \(E_b'\) its strong dual.
\end{enumerate}

\begin{enumerate}
\def\labelenumi{\alph{enumi})}
\item
  Show that every convex balanced subset \(V'\) of \(E_{b}'\) which absorbs the strongly bounded subsets of \(E_{b}'\) contains a barrel (for the strong topology) which absorbs the strongly bounded subsets of \(E_{b}'\) . (Let \((\mathbf{K}_{n}')\) be a countable base of the canonical bornology of \(E_{b}'\) and let \(\lambda_{n}\) be such that \(\lambda_{n}K_{n}' \subset \frac{1}{2}V'\) ; apply exerc. 2, b) to the sequence \(A_{n}'\) , where \(A_{n}'\) is the convex envelope of the union of the \(\lambda_{j}K_{j}'\) for \(j \leqslant n\) .)
\item
  Deduce from \(a\) ) that the following properties are equivalent :
\end{enumerate}

α) E is distinguished (IV, p.~52, exerc. 4).

β) \(E_{b}^{\prime}\) is infrabarrelled (III, p.~44, exerc. 7).

\(\gamma)\) \(\mathrm{E}_b^{\prime}\) is bornological.

δ) \(E_{b}^{\prime}\) is barrelled.

ε) \(E_{b}^{\prime}\) is ultrabornological (III, p.~45, exerc. 19).

\begin{enumerate}
\def\labelenumi{\alph{enumi})}
\setcounter{enumi}{2}
\tightlist
\item
  Show that if \(E_{b}^{\prime}\) is reflexive, then \(\hat{E} = E''\) (which is obviously reflexive) (cf.~IV, p.~52, exerc. 4 and p.~53, exerc. 11).
\end{enumerate}

\begin{enumerate}
\def\labelenumi{\arabic{enumi})}
\setcounter{enumi}{3}
\tightlist
\item
  Let E be a locally convex Hausdorff space, \(E'\) its dual. If M is a closed vector subspace of E which is metrizable and distinguished (IV, p.~52, exerc. 4), then the strong topology \(\beta(\mathrm{E}'/\mathrm{M}^{\circ}, \mathrm{M})\) is the quotient topology by \(M^{\circ}\) of the strong topology \(\beta(\mathrm{E}', \mathrm{E})\) (use exerc. 3, b) and IV, p.~51, exerc. 22, b)).
\end{enumerate}

¶ 5) For every integer n \textgreater{} 0, let \(a^{(n)}\) be the double sequence \((a_{pq}^{(n)})\) ( \(p \in N\) , \(q \in N\) ) such that \(a_{pq}^{(n)} = q\) if \(p \leqslant n\) and \(a_{pq}^{(n)} = 1\) if p \textgreater{} n.~Let E be the vector space of all double sequences \(x = (x_{pq})_{(p,q) \in \mathbf{N} \times \mathbf{N}}\) of real numbers such that, for every integer n \textgreater{} 0, the number \(r_n(x) = \sum_{p,q} a_{pq}^{(n)} |x_{pq}|\)

is finite. If E is assigned the topology defined by the semi-norms \(r_n\) , then E is a Fréchet space satisfying the first axiom of countability (IV, p.~47, exerc. 1, c)); the dual E' of E can be identified with the space of all double sequences \(x' = (x_{pq}')\) of real numbers such that for at least one index \(n\) , there exists \(k_n > 0\) such that \(|x_{pq}'| \leqslant k_n a_{pq}^{(n)}\) for every pair \((p, q)\) ; and \(\langle x, x' \rangle = \sum_{p, q} x_{pq} x_{pq}'\) .

(IV, p.~47, exerc. 1, c)).

For every integer \(p_0 > 0\) and every sequence \((m_p)\) of integers \(> 0\) , let \(\mathrm{J}(p_0; (m_p))\) be the set of pairs of integers \(p > 0\) , \(q > 0\) such that \(p \geqslant p_0\) and \(q \geqslant m_p\) ; let \(\mathfrak{V}\) be the filter base on \(\mathbf{N} \times \mathbf{N}\) consisting of the sets \(\mathrm{J}(p_0; (m_p))\) and let \(\mathfrak{F}\) be an ultrafilter which is finer than the filter with base \(\mathfrak{V}\) .

\begin{enumerate}
\def\labelenumi{\alph{enumi})}
\item
  Show that for all \(x' = (x_{pq}') \in \mathrm{E}'\) , the double sequence \((x_{pq}')\) has a limit \(u(x')\) with respect to the ultrafilter \(\mathfrak{F}\) ; if \(\mathrm{V}_n\) is a neighbourhood of 0 in \(\mathrm{E}\) defined by \(r_n(x) \leqslant 1\) , then \(|u(x')| \leqslant 1\) for all \(x' \in \mathrm{V}_n^\circ\) .
\item
  Let \(U'\) be a neighbourhood of 0 in \(E'\) , for the strong topology, which is convex, balanced and weakly closed, and for every n, let \(\alpha_{n} > 0\) be such that \(\alpha_{n} V_{n}^{\circ} \subset U'\) . For every integer p \textgreater{} 0, let \(m_{p}\) be an integer such that \(2^{p+1} \leqslant \alpha_{p} m_{p}\) , and let \(x' = (x'_{pq})\) be the double sequence with \(x'_{pq} = 0\) for q \textless{} \(m_{p}\) , \(x'_{pq} = 2\) for \(q \geqslant m_{p}\) . Show that \(x' \in U'\) but that \(u(x') = 2\) ; deduce that u is not strongly continuous in \(E'\) , while being bounded on every bounded subset of \(E'\) . Conclude (IV, p.~58, exerc. 3) that E is not distinguished, and consequently that the strong dual \(E'_{b}\) is a non infra-barrelled (DF) space.
\item
  Using b) construct an example of a closed subspace M of a Fréchet space F such that the strong topology \(\beta(\mathrm{F}^{\prime}/\mathrm{M}^{\circ}, \mathrm{M})\) is distinct from the quotient topology by \(M^{\circ}\) of the strong topology \(\beta(\mathrm{F}^{\prime}, \mathrm{F})\) (embed E in a countable product of Banach spaces).
\end{enumerate}

¶ 6) a) Let E be a (DF) space (IV, p.~57, exerc. 2), and let U be a convex set such that for every bounded, convex and balanced subset A of E, U ∩ A is a neighbourhood of 0 for the topology induced on A by that of E. Show that U is a neighbourhood of 0 in A. (Let (A \(_{n}\) ) be a countable base for the canonical bornology of E (III, p.~3, def. 5). Show that, by induction we can define a sequence ( \(\lambda_{n}\) ) of numbers \textgreater{} 0 and a sequence (V \(_{n}\) ) of closed, convex and balanced

neighbourhoods of 0 in E such that \(\lambda_{n}\mathrm{A}_{n} \subset \frac{1}{3}\mathrm{U}\) , \(\lambda_{n}\mathrm{A}_{n} \subset \bigcap_{j=1}^{\infty} \mathrm{V}_{j}\) , \(\mathrm{V}_{n} \cap \mathrm{A}_{n} \subset \mathrm{U}\) for every \(n\) .

First show that if \(\lambda_{j}\) and \(\mathrm{V}_j\) have been constructed for \(j\leqslant n\) , then we can find \(\lambda_{n + 1}\) such that \(\lambda_{n + 1}\mathrm{A}_{n + 1}\subset \frac{1}{3}\mathrm{U}\) and \(\lambda_{n + 1}\mathrm{A}_{n + 1}\subset \mathrm{V}_j\) for \(j\leqslant n\) . Next prove that we can find \(\mathrm{V}_{n + 1}\) such that \(\lambda_j\mathrm{A}_j\subset \mathrm{V}_{n + 1}\) for \(j\leqslant n + 1\) and \(\mathrm{V}_{n + 1}\cap \mathrm{A}_{n + 1}\subset \mathrm{U}\) ; for this, letting A denote the convex envelope of the \(\lambda_{j}\mathrm{A}_{j}\) for \(j\leqslant n + 1\) , show that we can take \(\mathrm{V}_{n + 1} = \overline{\mathrm{A} + \mathrm{V}}\) for a suitable convex balanced neighbourhood V of 0; we remark that for this it is enough to show that, if \(\mathbf{B} = \mathbf{A}_{n + 1}\cap \mathbb{C}\) U, then 0 is not in the closure of the set \(\mathbf{B} + 2\mathbf{A}\) .

\begin{enumerate}
\def\labelenumi{\alph{enumi})}
\setcounter{enumi}{1}
\tightlist
\item
  Deduce from a) that if u is a linear mapping from E into a locally convex space F, such that the restriction of u to every bounded subset of E is continuous, then u is continuous (cf.~IV, p.~50, exerc. 15).
\end{enumerate}

¶ 7) a) Let E be a (DF) space, U a convex, balanced and closed set in E, which absorbs the bounded subsets of E, and let \((x_n)\) be a sequence of points of \(\complement\) U. Show that there exists a neighbourhood V of 0 in E which does not contain any of the \(x_n\) . (Let \((A_n)\) be a countable base for the canonical bornology of E. Show that, by induction we can define a sequence \((\lambda_n)\) of numbers \textgreater{} 0 and a sequence \((V_n)\) of convex, balanced and closed neighbourhoods of

0 such that \(\lambda_{n}\mathbf{A}_{n}\subset \bigcap_{j = 1}\mathbf{V}_{j},\lambda_{n}\mathbf{A}_{n}\subset \mathrm{U}\) and \(x_{n}\in \mathbb{C}\mathbf{V}_{n}\) for all \(n\) . For this, if the \(\lambda_{j}\) and \(\mathrm{V}_j\) have

been constructed for \(j \leqslant n\) , take \(\lambda_{n+1}\) such that \(\lambda_{n+1}\mathrm{A}_{n+1} \subset \mathrm{U}\) and \(\lambda_{n+1}\mathrm{A}_{n+1} \subset \mathrm{V}_j\) for all \(j \leqslant n\) , then take \(\mathrm{V}_{n+1}\) containing the closure of the convex envelope of the union of the \(\lambda_j\mathrm{A}_j\) for \(j \leqslant n + 1\) .

\begin{enumerate}
\def\labelenumi{\alph{enumi})}
\setcounter{enumi}{1}
\item
  Deduce from a) that if M is a subset of E containing an everywhere dense countable set, then the topology induced on M by the strong topology of the bidual \(E''\) of E is identical with the topology induced by that of E. In particular, the convergent sequences in E are the same for the topology of E and for the topology induced by the strong topology of \(E''\) ; for every metrizable subset M of E, the topology induced on M by the topology of E is identical with the topology induced by the strong topology of \(E''\) .
\item
  Deduce from \(a\) ) that if there exists a countable everywhere dense set in \(E\) , then \(E\) is infra-barrelled.
\item
  Deduce from b) and from exerc. 6 that if every bounded subset of E is metrizable for the topology induced by that of E, then E is infrabarrelled.
\end{enumerate}

\begin{enumerate}
\def\labelenumi{\arabic{enumi})}
\setcounter{enumi}{7}
\tightlist
\item
  Let E be a Fréchet space, \(E_{b}^{\prime}\) its strong dual. Suppose that there exists an everywhere dense sequence \((x_{n}^{\prime})\) in \(E_{b}^{\prime}\) . Show that E satisfies the first axiom of countability. (Let \((\mathbf{K}_{n}^{\prime})\) be a countable base for the canonical bornology of \(E_{b}^{\prime}\) consisting of closed convex balanced sets. For every system \(\alpha\) consisting of a point \(x_{n}^{\prime}\) , an arbitrary finite number of rational numbers \(\lambda_{k} > 0 (1 \leqslant k \leqslant m)\) and m indices \(n_{k}\) such that \(x_{n}^{\prime} \notin 2 \sum_{k=1}^{m} \lambda_{k} K_{n_{k}}^{\prime} = 2 H_{\alpha}^{\prime}\) , let \(x_{\alpha} \in E\) be such that
\end{enumerate}

the hyperplane with equation \(\langle x_{\alpha}, y' \rangle = 1\) strictly separates the two weakly compact sets \(\mathrm{H}_{\alpha}'\) and \(x_{n}' + \mathrm{H}_{\alpha}'\) . Show that for every \(x' \neq 0\) in \(\mathrm{E}'\) , there exists a system \(\alpha\) such that \(\langle x_{\alpha}, x' \rangle \neq 0\) . For this, consider a neighbourhood \(\mathrm{V}'\) of 0 in \(\mathrm{E}_b'\) such that \(\mathrm{V}' \cap (x' + \mathrm{V}') = \emptyset\) , then for each integer \(m\) , take a rational number \(\lambda_m > 0\) such that \(\lambda_m \mathrm{K}_m' \subset \mathrm{V}'\) ; use the fact that the union \(\mathrm{U}' \subset \mathrm{V}'\) of the \(\lambda_m \mathrm{K}_m'\) is a neighbourhood of 0 (exerc. 7, and IV, p.~58, exerc. 3, b)) and that there exists \(n\) such that \(x_n' \in x' + \mathrm{U}'\) .

¶ 9) a) Let E be a Hausdorff semi-barrelled space, M a closed vector subspace of E and E' the dual of E. Show that E/M is semi-barrelled and that the strong topology β(M°, E/M) is identical with the topology induced on M° by the strong topology β(E', E). (Note that it is enough to prove that a sequence (x') in M° which converges to 0 for β(E', E) is bounded for β(M°, E/M).) Deduce that if E is a (DF) space, then so is E/M (cf.~IV, p.~63, exerc. 8).

\begin{enumerate}
\def\labelenumi{\alph{enumi})}
\setcounter{enumi}{1}
\item
  Let E be a locally convex Hausdorff space, M a (not-necessarily closed) vector subspace of E. Show that if M is a semi-barrelled space, then the strong topology \(\beta(\mathrm{E}^{\prime}/\mathrm{M}^{\circ}, \mathrm{M})\) is identical with the quotient topology by \(M^{\circ}\) of the strong topology \(\beta(\mathrm{E}^{\prime}, \mathrm{E})\) . (Argue as in a).)
\item
  Show that a Hausdorff and quasi-complete semi-barrelled space E is complete (use b) applied to E and \(\hat{E}\) ). In particular, a semi-barrelled, semi-reflexive space is complete (cf.~IV, p.~52, exerc. 6).
\item
  Show that the completion of a semi-barrelled (resp. (DF)) Hausdorff space is semi-barrelled (resp. a (DF) space).
\item
  Let \((\mathrm{E}_n)\) be a sequence of semi-barrelled (resp. (DF)) spaces, E a vector space, and for each \(n\) , let \(f_n\) be a linear mapping from \(\mathrm{E}_n\) into E. Suppose that \(\dot{\mathrm{E}}\) is the union of the \(f_n(\mathrm{E}_n)\) ; show that, E is semi-barrelled (resp. a (DF) space) for the finest locally convex topology for which all the \(f_n\) are continuous (first examine the case where E is the topological direct sum of the \(\mathrm{E}_n\) ). If the \(\mathrm{E}_n\) are semi-reflexive (resp. reflexive) and if E is Hausdorff, then E is semi-reflexive (resp. reflexive).
\end{enumerate}

\begin{enumerate}
\def\labelenumi{\arabic{enumi})}
\setcounter{enumi}{9}
\tightlist
\item
  Let E be a Fréchet space satisfying the first axiom of countability. Show that if in the dual \(E'\) of E, every sequence which converges for the weak topology \(\sigma(E', E)\) also converges for the strong topology \(\beta(E', E)\) , then E is a Montel space. (Show that every bounded subset of \(E'\) is relatively compact for the strong topology; use GT, II, § 4, exerc. 6; then use IV, p.~53, exerc. 11, b).)
\end{enumerate}

¶ 11) Let \((c_{mn})\) be a double sequence of numbers \textgreater{} 0 such that \(c_{m,n} \leqslant c_{m+1,n}\) and let E be the space of all sequences \(x = (x_n)\) of real numbers such that \(p_m(x) = \sum_n c_{mn} |x_n| < +\infty\) for

every integer \(m\) . We endow \(E\) with the topology defined by the semi-norms \(p_m\) and for this topology \(E\) is a Fréchet space satisfying the first axiom of countability; the dual \(E'\) of \(E\) can be identified with the space of all sequences \(x' = (x_n')\) such that \(\sup_{n} c_{mn}^{-1} |x_n'| < +\infty\) for at

least one \(m\) , the canonical bilinear form \(\langle x, x' \rangle\) being identified with \(\sum_{n} x_n x_n'\) (IV, p.~47,

exerc. 1, c)). Suppose that there does not exist any subsequence \((n_k)\) for which there is a sequence \((a_m)\) of numbers \(\geqslant 0\) and an index \(m_0\) such that \(c_{m,n_k} \leqslant a_m c_{m_0,n_k}\) for all \(m \geqslant m_0\) and for all \(k\) . Under these conditions, every weakly convergent sequence in \(\mathbf{E}'\) is strongly convergent and consequently (IV, p.~60, exerc. 10) E is a Montel space. (Argue by reductio ad absurdum; if necessary make a transformation of the form \((x_n) \mapsto (a_n x_n)\) to reduce to the case where \(c_{m_0 n} = 1\) for all \(n\) and some \(m_0\) , and where there exists a sequence \((x'^{(p)})_{p \geqslant 0}\) which converges weakly to 0 in \(\mathbf{E}'\) , and is such that \(|x_n'^{(p)}| \leqslant 1\) for every pair \((p, n)\) , and also that there exists a bounded set B in E, defined by \(p_m(x) \leqslant b_m\) for all \(m \geqslant 0\) and such that \(\sup_{x \in \mathbf{B}} |\langle x, x'^{(p)} \rangle| \geqslant 2\delta > 0\) for

every integer \(p\) . Under these hypothesis, prove that there exists a strictly increasing sequence \((r_q)\) of integers, and a sequence \((x^{(q)})\) of points of \(\mathbf{B}\) such that

\[
\sum_ {k = r _ {q} + 1} ^ {r _ {q+1}} \left| x _ {k} ^ {(q)} \right| > \delta
\]

for each index \(q\) . Then show, by reductio ad absurdum, that for every \(q\) , there exists at least one index \(s_q\) such that \(r_q < s_q \leqslant r_{q+1}\) and that for every integer \(m\) , we have \(c_{m,s_q} \leqslant b_m 2^{m+2}/\delta\) , which contradicts the hypothesis.)

\begin{enumerate}
\def\labelenumi{\arabic{enumi})}
\setcounter{enumi}{11}
\tightlist
\item
  \begin{enumerate}
  \def\labelenumii{\alph{enumii})}
  \tightlist
  \item
    Let F be a Hausdorff (DF) space, and \(F_{b}^{\prime}\) its strong dual. Show that if \(F_{b}^{\prime}\) is reflexive, then the completion \(\hat{F}\) of F is reflexive and is equal to the bidual \(F^{\prime\prime}\) of F (cf.~IV, p.~52, exerc. 4 and p.~53, exerc. 11).
  \end{enumerate}
\end{enumerate}

\begin{enumerate}
\def\labelenumi{\alph{enumi})}
\setcounter{enumi}{1}
\tightlist
\item
  Let E be a Fréchet space. Show that if the bidual \(E''\) of E is reflexive, then E is reflexive.
\end{enumerate}

\begin{enumerate}
\def\labelenumi{\arabic{enumi})}
\setcounter{enumi}{12}
\tightlist
\item
  \begin{enumerate}
  \def\labelenumii{\alph{enumii})}
  \tightlist
  \item
    Let E, F be two Fréchet spaces, G a locally convex Hausdorff space and \(E'\) , \(F'\) , \(G'\) the duals of E, F, G respectively. Let u be a bilinear mapping from \(E' \times F'\) into \(G'\) , which is separately continuous (III, p.~28) when \(E'\) , \(F'\) , \(G'\) are assigned the weak topologies \(\sigma(E', E)\) , \(\sigma(F', F)\) and \(\sigma(G', G)\) . Show that under these conditions, u is a continuous mapping from \(E' \times F'\) into \(G'\) when \(E'\) , \(F'\) and \(G'\) are assigned the strong topologies. (For \(z \in G\) , put \(\langle z, u(x', y') \rangle = \langle v_z(x'), y' \rangle\) where \(v_z(x') \in F\) . First show that if \(E'\) is assigned the strong topology and \(F\) the initial topology, then the set of all \(v_z\) , as \(z\) ranges over a bounded set \(C\) in \(G\) , is equicontinuous; for this use IV, p.~51, exerc. 19, d). Next show that there exists a neighbourhood \(V'\) of 0 for the strong topology of \(E'\) such that the union of the sets \(v_z(V')\) as \(z\) ranges over \(C\) is bounded in \(F\) ; for this use III, p.~47, exerc. 5.)
  \end{enumerate}
\end{enumerate}

\begin{enumerate}
\def\labelenumi{\alph{enumi})}
\setcounter{enumi}{1}
\tightlist
\item
  Give an example to show that the conclusion of a) does not hold if we assume that E is a Fréchet space and F a strict inductive limit of Fréchet spaces (III, p.~47, exerc. 3).
\end{enumerate}

\begin{enumerate}
\def\labelenumi{\arabic{enumi})}
\setcounter{enumi}{13}
\tightlist
\item
  \begin{enumerate}
  \def\labelenumii{\alph{enumii})}
  \tightlist
  \item
    Let E be a Fréchet space, \(E'\) its dual. Show that \(E'\) , endowed with the topology of compact convergence or with a finer S-topology, is complete (cf.~III, p.~22, Remark 1). If E is not reflexive, show that \(E'\) is not infrabarrelled for any S-topology which is finer than the topology of compact convergence and coarser than \(\tau(\mathrm{E}', \mathrm{E})\) .
  \end{enumerate}
\end{enumerate}

\begin{enumerate}
\def\labelenumi{\alph{enumi})}
\setcounter{enumi}{1}
\tightlist
\item
  Let \((\mathrm{E}_{\alpha})_{\alpha\in\mathbf{A}}\) be a family of Fréchet spaces, E a vector space and for every \(\alpha\in A\) , let \(h_{\alpha}\) be a linear mapping from \(E_{\alpha}\) into E. Suppose that E, endowed with the finest locally convex topology for which the \(h_{\alpha}\) are continuous (II, p.~27) is Hausdorff. Prove that the dual \(E'\) of E, endowed with the topology of compact convergence or with any finer S-topology, is complete (cf.~III, p.~20, th. 1).
\end{enumerate}

\begin{enumerate}
\def\labelenumi{\arabic{enumi})}
\setcounter{enumi}{14}
\tightlist
\item
  Let \(\mathrm{E}\) be an infinite dimensional Banach space, \((a_{n})_{n\geqslant 1}\) a countable total free family of points of \(\mathrm{E}\) , and let \(\mathrm{F}_n\) be the \(n\) -dimensional subspace of \(\mathrm{E}\) generated by the \(a_{m}\) for \(m\leqslant n\) . Let \(\mathrm{S}_n\) be the sphere with equation \(\| x\| = n\) in \(\mathrm{E}\) ; in \(\mathrm{S}_n\cap \mathrm{F}_n\) let \(\mathrm{A}_n\) be a finite set such that every
\end{enumerate}

point of \(S_{n} \cap F_{n}\) is at a distance \(\leqslant 1/n\) from \(A_{n}\) . Prove that \(A = \bigcap_{n=1}^{\infty} A_{n}\) is such that its inter-

section with every closed bounded set is closed, but that 0 is a limit point of A for the weakened topology.

\begin{enumerate}
\def\labelenumi{\arabic{enumi})}
\setcounter{enumi}{15}
\tightlist
\item
  Let E be an inductive limit space of a sequence \((\mathrm{E}_{p})\) of locally convex metrizable spaces, the canonical mappings \(E_{p} \rightarrow E\) being injective and let E be the union of the images of the \(E_{p}\) . Show that the strong dual \(E_{b}^{\prime}\) of E is exhaustible (III, p.~49, exerc. 1). (Let \((\mathrm{U}_{j}^{p})\) be a decreasing fundamental system of closed convex and balanced neighbourhoods of 0 in \(E_{p}\) ; consider the finite intersections of the polar sets \((\mathrm{U}_{j}^{p})^{\circ}\) in \(E^{\prime}\) ; use the fact that for every increasing sequence
\end{enumerate}

\((m_{j})_{j\geqslant 1}\) the intersection \(\bigcap_{j = 1}^{\infty}(\mathrm{U}_{m_j}^j)^{\circ}\) is the polar of a neighbourhood of 0 in E, and that \(\mathrm{E}_b^\prime\) is complete.)

¶ 17) a) Let E be a locally convex Hausdorff space, such that the bornology consisting of the relatively compact sets of E has a countable base (A \(_{n}\) ) (III, p.~1). Show that (A \(_{n}\) ) is also a base for the canonical bornology (III, p.~3, def. 5). (Let C \(_{n}\) be the relatively compact set which is the sum of n sets each equal to A \(_{n}\) ; then (C \(_{n}\) ) is also a base for the bornology of relatively compact sets. Argue by reductio ad absurdum, considering a bounded set B which is not contained in any of the C \(_{n}\) , and conclude that there exists a sequence (x \(_{n}\) ) of points of B such that x \(_{n}\) /n∉ A \(_{n}\) ; deduce a contradiction.) Then, the space E is semi-reflexive and the closed convex balanced envelope of every compact set in E is compact.

\begin{enumerate}
\def\labelenumi{\alph{enumi})}
\setcounter{enumi}{1}
\tightlist
\item
  Suppose that E is infrabarrelled and that for a topology T compatible with the duality between E and \(E'\) , there exists a countable base for the bornology of relatively compact subsets of E for T. Show that then E is a reflexive (DF) space (IV, p.~57, exerc. 2). (Use a), observing that the closed bounded sets in E are compact for T, and consequently, complete for the initial topology of E; this implies that E is barrelled). If T is the initial topology, then E is a Montel space.
\end{enumerate}

\begin{enumerate}
\def\labelenumi{\arabic{enumi})}
\setcounter{enumi}{17}
\tightlist
\item
  \begin{enumerate}
  \def\labelenumii{\alph{enumii})}
  \tightlist
  \item
    Let E be a locally convex Hausdorff space, and let \((\mathbf{A}_{n})\) be an increasing sequence of convex, balanced, compact sets for the weakened topology \(\sigma(\mathrm{E}, \mathrm{E}')\) , such that the union of the sets \(A_{n}\) is E, and that for every integer n and every \(\lambda > 0\) , there exists m such that \(\lambda A_{n} \subset A_{m}\) . Show that \((\mathbf{A}_{n})\) is a base for the bornology consisting of convex and relatively compact sets for \(\sigma(\mathrm{E}, \mathrm{E}')\) . (Observe that on \(E'\) the topology of uniform convergence on the \(A_{n}\) is \(\tau(\mathrm{E}', \mathrm{E})\) .) b) If E is barrelled and if there exists a sequence \((\mathbf{A}_{n})\) having the properties mentioned in a), then E is a reflexive (DF) space. (Observe that \(E'\) , endowed with \(\tau(\mathrm{E}', \mathrm{E})\) is metrizable and that the topology of E is \(\beta(\mathrm{E}, \mathrm{E}')\) .)
  \end{enumerate}
\end{enumerate}

